\renewcommand{\bibname}{Quellen und Bandzuordnung}
\begin{thebibliography}{9}
\addcontentsline{toc}{chapter}{Quellen und Bandzuordnung}
\footnotesize
\setlength{\itemsep}{6pt}
\setlength{\parskip}{0pt}
\interlinepenalty=10000

\bibitem{project}
Martin Kunze: \emph{Die Grundlagen der Mathematik}.
Band~14: Existenzprädikate für Paarwerte; universelle Eigenschaften;
Idempotenz, Kommutativität und Assoziativität; Paarwerte in
Mengenfamilien und im Potenzmengenträger.
Kapitel~\ref{b14guide:ch:foundations},
\ref{b14guide:ch:characterizations},
\ref{b14guide:ch:iteration} und \ref{b14guide:ch:sets}
fassen diese Aussagen zusammen.
Die übrigen Regeln sind unmittelbare Folgerungen aus den
Paarbegriffen und ihren universellen Eigenschaften.
Format nach den Leitfäden zu den Bänden~3, 4, 5, 28 und~37.

\bibitem{bounds}
The mathlib Community: \emph{Mathlib.Order.Bounds.Defs}
und \emph{Mathlib.Order.Bounds.Basic}, offizielle Dokumentation,
abgerufen am 23.~Oktober 2026.
Abgleich der Paarbegriffe über \path{IsGLB} und \path{IsLUB};
der Kandidatenkriterien über \path{isGLB_iff_le_iff} und
\path{isLUB_iff_le_iff}; der Eindeutigkeit über
\path{IsGLB.unique} und \path{IsLUB.unique}.
Hier werden diese Kriterien ausschließlich auf $\{x,y\}$ angewandt.
\url{https://leanprover-community.github.io/mathlib4_docs/Mathlib/Order/Bounds/Defs.html}.
\url{https://leanprover-community.github.io/mathlib4_docs/Mathlib/Order/Bounds/Basic.html}.

\bibitem{lattice}
The mathlib Community: \emph{Mathlib.Order.Lattice},
offizielle Dokumentation, abgerufen am 23.~Oktober 2026.
Vergleichsstellen für Paaridentitäten und strikte Vergleiche:
\path{inf_lt_left}, \path{left_lt_sup}, \path{inf_eq_sup},
\path{inf_lt_sup}, \path{sup_eq_sup_mono_left}
und die dualen Aussagen.
\url{https://leanprover-community.github.io/mathlib4_docs/Mathlib/Order/Lattice.html}.

\end{thebibliography}

\begingroup
\footnotesize
\noindent\textbf{Umfang des Quellenabgleichs.}
Die Bibliothek verwendet für viele Paaridentitäten überall
definierte Operationen und stärkere Strukturklassen.
Dieser Leitfaden übernimmt daraus keine Existenzaxiome:
Seine Formulierungen benutzen nur die ausdrücklich genannten
lokalen Paarwerte. Die lokalen Existenzverschärfungen, die
Regeln für zulässige Paarterme und die gemischten Ungleichungen
sind hier aus den Paarbedingungen formulierte Folgerungen;
die Quellen werden nicht als wortgleicher Nachweis jeder
Formulierung angeführt.
Es werden weder neue Forschungsergebnisse noch historische
Erstentdeckungen beansprucht.
\endgroup